\documentclass{ip-quiz}

\quizlevel{n1}
\quiznum{2}
\quiztitle{Funciones - Opción 1}

\begin{document}

\makeheader
\maketitle

\begin{sectionbox}{Enunciado}
Considere un cuadrado de lado $x$ inscrito en un círculo. Los vértices del cuadrado tocan la circunferencia.

\begin{center}
\begin{tikzpicture}[scale=0.8]
    % Circle
    \draw[thick] (2,2) circle (2);
    
    % Square inscribed (rotated 45 degrees, vertices on circle)
    \draw[thick] (2,0) -- (4,2) -- (2,4) -- (0,2) -- cycle;
    
    % Side label
    \draw[{Stealth}-{Stealth}] (2,0) -- (4,2) node[midway, above left] {$x$};
\end{tikzpicture}
\end{center}

Defina una función llamada \texttt{calcular\_area\_circulo} que reciba como parámetro el lado $x$ del cuadrado y retorne el área del círculo. Aproxime $\pi$ como \texttt{3.14}.
\end{sectionbox}

\begin{sectionbox}{Solución}
\begin{minted}{python}
def calcular_area_circulo(x: float) -> float:
    cateto_al_cuadrado = x ** 2
    hipotenusa_al_cuadrado = cateto_al_cuadrado * 2
    diagonal_del_cuadrado = hipotenusa_al_cuadrado ** 0.5
    radio_circulo = diagonal_del_cuadrado / 2
    area_circulo = 3.14 * radio_circulo ** 2
    return area_circulo
\end{minted}
\end{sectionbox}
    
\end{document}